\documentclass[11pt,a4paper]{article}
\usepackage[margin=23mm,headheight=14pt,headsep=8.5mm,footskip=10mm]{geometry}
\usepackage{amsmath,amssymb}
\usepackage[no-math]{fontspec}
\usepackage{xeCJK}
\setmainfont{Latin Modern Roman}
\setCJKmainfont{FandolSong-Regular.otf}[Path=assets/fonts/fandol/,BoldFont=FandolSong-Bold.otf,ItalicFont=FandolSong-Regular.otf]
\setCJKsansfont{FandolSong-Regular.otf}[Path=assets/fonts/fandol/,BoldFont=FandolSong-Bold.otf]
\xeCJKsetup{PunctStyle=quanjiao,CheckSingle=true}
\usepackage{setspace}\setstretch{1.12}
\setlength{\parindent}{2em}\setlength{\parskip}{3pt}
\usepackage{titlesec}
\titleformat{\section}{\centering\normalfont\Large\bfseries}{}{0pt}{}
\titlespacing*{\section}{0pt}{1.5\baselineskip}{0.8\baselineskip}
\usepackage{needspace,enumitem,longtable,booktabs,array,calc}
\providecommand{\tightlist}{\setlength{\itemsep}{3pt}\setlength{\parskip}{0pt}}
\usepackage{fancyhdr}\pagestyle{fancy}\fancyhf{}
\fancyhead[L]{\small 数学手记}\fancyhead[R]{\small 楼红卫 B 组 · 第四十六批}
\fancyfoot[C]{\thepage}\renewcommand{\headrulewidth}{0.4pt}
\usepackage[hidelinks,unicode,linktoc=all]{hyperref}\usepackage{bookmark}
\hypersetup{pdftitle={楼红卫数学分析 B317—B322：Fourier 变换的延拓与级数求和},pdfauthor={Vincent Arqain}}
\renewcommand{\contentsname}{题号目录}\setcounter{tocdepth}{1}\setcounter{secnumdepth}{0}
\widowpenalty=10000\clubpenalty=10000\displaywidowpenalty=10000
\predisplaypenalty=10000\interdisplaylinepenalty=10000
\raggedbottom\emergencystretch=2em
\begin{document}
\begin{center}
{\LARGE\bfseries 楼红卫《数学分析》B 组习题解答\par}
\vspace{7pt}{\large 第四十六批\quad B317—B322\par}
\vspace{5pt}{Fourier 变换的延拓与级数求和\par}
\end{center}
习题 12.3.B 第 5—9 题及 12.4.B 第 1 题。本文采用指数 \(e^{-2\pi i x\cdot\xi}\) 的 Fourier 变换。
\begingroup\setlength{\parskip}{0pt}\tableofcontents\endgroup

\Needspace{8\baselineskip}
\section{问题 B317}\label{ux95eeux9898-b317}

\textbf{原题号：} 习题 12.3.B 第 5 题。原书下册第 386
页。\href{https://zj139.github.io/math-notes/files/problems/楼分析习题-v2.pdf\#page=43}{题册
PDF 第 43 页}

\textbf{题目}　设 \(1<p<2\)，\(p'=p/(p-1)\)。若速降函数列 \(\varphi_k\)
在 \(L^p(\mathbb R^n)\) 中强收敛于 \(f\)，证明 \(\widehat{\varphi_k}\)
在 \(L^{p'}(\mathbb R^n)\) 中强收敛于 \(\widehat f\)。

\textbf{证明}　先在速降函数空间 \(\mathcal S\) 使用原书定理 12.3.8。将
Hausdorff--Young 不等式用于两个逼近函数之差，得到

\begin{equation}
\|\widehat{\varphi_k}-\widehat{\varphi_j}\|_{p'}
\leq \|\varphi_k-\varphi_j\|_p .\tag{1}
\end{equation}

右端随 \(j,k\to\infty\) 趋于零。由 \(L^{p'}\) 完备，左侧的函数列有强极限
\(g\)。若另一列 \(\psi_k\in\mathcal S\) 也趋于 \(f\)，则其变换与
\(\widehat{\varphi_k}\) 的距离至多为
\(\|\varphi_k-f\|_p+\|f-\psi_k\|_p\)，故极限相同。

还须说明这个极限与教材在 \(L^1+L^2\) 上定义的变换相同。任取
\(\eta\in\mathcal S\)，先对 \(\varphi_k\) 用 Fubini 定理，再用两次
Hölder 不等式传递极限，得到

\begin{equation}
\int_{\mathbb R^n} g(\xi)\eta(\xi)\,d\xi
=\lim_{k\to\infty}\int_{\mathbb R^n}\varphi_k(x)\widehat\eta(x)\,dx
=\int_{\mathbb R^n} f(x)\widehat\eta(x)\,dx .\tag{2}
\end{equation}

这里采用双线性配对，故两边均没有共轭。写 \(f=u+v\)，其中
\(u=f\mathbf1_{\{|f|>1\}}\)，\(v=f\mathbf1_{\{|f|\le1\}}\)。由 \(1<p<2\)
可知 \(u\in L^1\)、\(v\in L^2\)。对 \(u\) 直接用 Fubini 定理，对 \(v\)
用速降函数作 \(L^2\) 逼近并用速降类的 Plancherel 等式，便知
\(\widehat u+\widehat v\)
也满足（2）。两者都是局部可积函数，对所有光滑紧支测试函数的积分相同，故几乎处处相等。这就识别了
\(g=\widehat f\)。在（1）中固定 \(k\) 令 \(j\to\infty\)，还得到误差界
\(\|\widehat{\varphi_k}-\widehat f\|_{p'}\leq\|\varphi_k-f\|_p\)。

\Needspace{8\baselineskip}
\section{问题 B318}\label{ux95eeux9898-b318}

\textbf{原题号：} 习题 12.3.B 第 6 题。原书下册第 386
页。\href{https://zj139.github.io/math-notes/files/problems/楼分析习题-v2.pdf\#page=43}{题册
PDF 第 43 页}

\textbf{题目}　设
\(f\in L^1(\mathbb R^n)\cap L^2(\mathbb R^n)\)。构造同一列速降函数
\(f_k\)，使它在两种范数下都强收敛于 \(f\)。

\textbf{证明}　先截去远处和大函数值。令
\(u_k=f\mathbf1_{\{|x|\le k,\ |f(x)|\le k\}}\)。对 \(r=1,2\)，由
\(|u_k-f|^r\le |f|^r\) 及支配收敛定理，\(\|u_k-f\|_r\to0\)。此时 \(u_k\)
已有紧支撑，但还需光滑化。

取非负 \(\rho\in C_c^\infty(\mathbb R^n)\)，支撑在单位球中且积分为一，置
\(\rho_\varepsilon(x)=\varepsilon^{-n}\rho(x/\varepsilon)\)。记
\(\tau_yu(x)=u(x-y)\)。由 Minkowski 不等式，

\begin{equation}
\|\rho_\varepsilon*u_k-u_k\|_r
\leq\int \rho(z)\|\tau_{\varepsilon z}u_k-u_k\|_r\,dz,
\qquad r=1,2 .\tag{3}
\end{equation}

对固定
\(k\)，平移在这两种范数下都连续。事实上，对连续紧支函数，这由一致连续性和支撑有界可知；再用连续紧支函数在
\(L^r\) 中的稠密性及平移等距即得一般情形。式（3）被
\(2\rho(z)\|u_k\|_r\) 控制，故右端随 \(\varepsilon\downarrow0\)
趋于零。对每个 \(k\) 选同一个 \(0<\varepsilon_k<1/k\)，使两种误差都小于
\(1/k\)。于是

\begin{equation}
f_k=\rho_{\varepsilon_k}*u_k\in C_c^\infty(\mathbb R^n)\subset\mathcal S,
\qquad \|f_k-f\|_r\leq k^{-1}+\|u_k-f\|_r\longrightarrow0 .\tag{4}
\end{equation}

卷积的各阶导数由对 \(\rho_{\varepsilon_k}\) 求导给出，支撑包含在半径
\(k+\varepsilon_k\) 的球内。这也验证了所构造函数确属速降类。

\Needspace{8\baselineskip}
\section{问题 B319}\label{ux95eeux9898-b319}

\textbf{原题号：} 习题 12.3.B 第 7 题。原书下册第 386
页。\href{https://zj139.github.io/math-notes/files/problems/楼分析习题-v2.pdf\#page=43}{题册
PDF 第 43 页}

\textbf{题目}　证明 Plancherel 定理在 \(L^2(\mathbb R^n)\) 中成立。

\textbf{证明}　对 \(f\in L^2\)，取 \(\varphi_k\in\mathcal S\) 在 \(L^2\)
中逼近 \(f\)。原书定理 12.3.6 给出速降函数的等距性，故

\begin{equation}
\|\widehat{\varphi_k}-\widehat{\varphi_j}\|_2
=\|\varphi_k-\varphi_j\|_2 .\tag{5}
\end{equation}

由完备性，\(\widehat{\varphi_k}\) 有 \(L^2\) 极限，记为
\(\mathcal F_2 f\)。对两列逼近函数之差再用（5），可知定义与选列无关。线性由逐项取极限得到，范数连续性给出

\begin{equation}
\int_{\mathbb R^n}|\mathcal F_2 f(\xi)|^2\,d\xi
=\lim_{k\to\infty}\|\widehat{\varphi_k}\|_2^2
=\lim_{k\to\infty}\|\varphi_k\|_2^2
=\int_{\mathbb R^n}|f(x)|^2\,dx .\tag{6}
\end{equation}

其像是闭的。确实，像中的收敛列由等距性拉回为 Cauchy
列，原像的极限又映到给定极限。另一方面，速降函数的反演公式表明
\(\mathcal F(\mathcal S)=\mathcal S\)，故像还包含稠密子集。由此
\(\mathcal F_2\) 是 \(L^2\) 到自身的满等距线性映射，逆算子是逆 Fourier
变换的连续延拓。

若 \(f\in L^1\cap L^2\)，用 B318 选取同时逼近的一列。此时
\(\|\widehat{\varphi_k}-\widehat f_{\!1}\|_\infty\leq\|\varphi_k-f\|_1\to0\)，其中
\(\widehat f_{\!1}\) 由绝对收敛积分定义。同时
\(\widehat{\varphi_k}\to\mathcal F_2 f\) 于
\(L^2\)，抽取子列后几乎处处收敛。因此两种定义几乎处处相同。对一般
\(L^2\) 函数，式（6）中的变换由均方极限给出。

\Needspace{8\baselineskip}
\section{问题 B320}\label{ux95eeux9898-b320}

\textbf{原题号：} 习题 12.3.B 第 8 题。原书下册第 386
页。\href{https://zj139.github.io/math-notes/files/problems/楼分析习题-v2.pdf\#page=43}{题册
PDF 第 43 页}

\textbf{题目}　设 \(1\leq p\leq2\)。证明对于每个
\(f\in L^p(\mathbb R^n)\) 有 Hausdorff--Young 不等式

\begin{equation}
\|\widehat f\|_{p'}\leq\|f\|_p,
\qquad \frac{1}{p}+\frac{1}{p'}=1 .\tag{7}
\end{equation}

\textbf{证明}　当 \(p=1\)，变换积分的指数因子模为一，逐点取绝对值得
\(|\widehat f(\xi)|\leq\int|f|\)。当 \(p=2\)，式（7）就是 B319
中的等式（6）。

设 \(1<p<2\)。原书定理 12.3.8 已在 \(\mathcal S\) 上给出（7）。取
\(\varphi_k\in\mathcal S\) 使 \(\varphi_k\to f\) 于 \(L^p\)。由
B317，\(\widehat{\varphi_k}\to\widehat f\) 于 \(L^{p'}\)。因此

\begin{equation}
\|\widehat f\|_{p'}
=\lim_{k\to\infty}\|\widehat{\varphi_k}\|_{p'}
\leq\lim_{k\to\infty}\|\varphi_k\|_p
=\|f\|_p .\tag{8}
\end{equation}

B317 用的是速降类上已证明的不等式。此处借助强收敛把它传给任意 \(L^p\)
函数，全程只涉及范数极限。

\Needspace{8\baselineskip}
\section{问题 B321}\label{ux95eeux9898-b321}

\textbf{原题号：} 习题 12.3.B 第 9 题。原书下册第 386
页。\href{https://zj139.github.io/math-notes/files/problems/楼分析习题-v2.pdf\#page=43}{题册
PDF 第 43 页}

\textbf{题目}　推广定理 12.3.4 和定理
12.3.5。原书在速降函数上分别给出卷积与乘积的变换公式，以及保内积公式。

\textbf{解}　以下三式将原书的两个定理分别推广到可积函数与平方可积函数。仍用
\(\widehat f(\xi)=\int f(x)e^{-2\pi i x\cdot\xi}\,dx\) 的归一化。

\textbf{卷积。} 若 \(f,g\in L^1\)，则 \(f*g\in L^1\)，且
\(\widehat{f*g}(\xi)=\widehat f(\xi)\widehat g(\xi)\) 对每个 \(\xi\)
成立。事实上，由 Tonelli 定理，

\begin{equation}
\iint |f(y)g(x-y)|\,dx\,dy=\|f\|_1\|g\|_1 .\tag{9}
\end{equation}

对变换的二重积分用 Fubini 定理，令 \(z=x-y\)，指数因子分解为关于 \(y,z\)
的两个因子，便得到结论。

\textbf{乘积。} 若 \(f,g\in L^2\)，则 \(fg\in L^1\)，并且对每个 \(\xi\)
有

\begin{equation}
\widehat{fg}(\xi)=
\int_{\mathbb R^n}\widehat f(\eta)\widehat g(\xi-\eta)\,d\eta .\tag{10}
\end{equation}

右端对每个 \(\xi\) 都绝对收敛，由 Cauchy--Schwarz
不等式，其绝对积分至多为 \(\|f\|_2\|g\|_2\)。取 \(f_k,g_k\in\mathcal S\)
分别在 \(L^2\) 中逼近 \(f,g\)。由

\begin{equation}
\|f_kg_k-fg\|_1\leq
\|f_k-f\|_2\|g_k\|_2+\|f\|_2\|g_k-g\|_2\longrightarrow0,\tag{11}
\end{equation}

左端变换一致收敛。右端卷积之差拆为
\((\widehat f_k-\widehat f)*\widehat g_k+\widehat f*(\widehat g_k-\widehat g)\)。逐项用
Cauchy--Schwarz 和 B319，同样得到（11）右端的界，且与 \(\xi\)
无关。对速降函数的公式取一致极限即得（10）。等式两边给出连续函数，改变
\(L^2\) 代表元也不影响任何一点的积分。

\Needspace{4\baselineskip}\textbf{内积。} 若 \(f,g\in L^2\)，则

\begin{equation}
\int_{\mathbb R^n} f(x)\overline{g(x)}\,dx
=\int_{\mathbb R^n}\widehat f(\xi)\overline{\widehat g(\xi)}\,d\xi .\tag{12}
\end{equation}

两积分都由 Cauchy--Schwarz 保证绝对收敛。对上述 \(f_k,g_k\) 用原书定理
12.3.5。原空间的内积误差由（11）右端控制，频率空间的内积误差由同一上界控制，令
\(k\to\infty\) 即得（12）。特别取 \(g=f\) 就恢复 Plancherel 等式。

\Needspace{8\baselineskip}
\section{问题 B322}\label{ux95eeux9898-b322}

\textbf{原题号：} 习题 12.4.B 第 1 题。原书下册第 393
页。\href{https://zj139.github.io/math-notes/files/problems/楼分析习题-v2.pdf\#page=43}{题册
PDF 第 43 页}

\textbf{题目}　推广引理 12.4.5。原引理断言，若
\(\sum_{n\ge1}c_n=A\)，则以 \((\sin(nh)/(nh))^2\) 加权后求和，在
\(h\to0\) 时仍趋于 \(A\)。

\textbf{解}　原证明中能统一控制求和误差的是权函数的总变差。设
\(K:[0,\infty)\to\mathbb C\) 有限变差，在零点右连续，\(K(0)=1\)，且
\(\lim_{t\to\infty}K(t)=0\)。则对每个收敛级数 \(\sum c_n=A\)，有

\begin{equation}
\lim_{h\to0}\sum_{n=1}^\infty c_nK(n|h|)=A .\tag{13}
\end{equation}

对每个固定 \(h\ne0\)，左边的级数也收敛。令
\(S_n=\sum_{j=1}^n c_j\)，\(S_0=0\)，\(w_n=K(n|h|)\)。对前 \(N\) 项作
Abel 分部求和，得

\begin{equation}
\sum_{n=1}^N c_n w_n
=S_N w_N+\sum_{n=1}^{N-1}S_n(w_n-w_{n+1}) .\tag{14}
\end{equation}

记 \(V\) 为 \(K\) 的总变差。因
\(\sum_{n\ge1}|w_n-w_{n+1}|\leq V\)，\(S_n\) 有界且
\(w_N\to0\)，（14）右边的极限存在。又
\(\sum_{n\ge1}(w_n-w_{n+1})=w_1\)，从而

\begin{equation}
\sum_{n\ge1} c_nw_n-A
=A(w_1-1)+\sum_{n\ge1}(S_n-A)(w_n-w_{n+1}) .\tag{15}
\end{equation}

给定 \(\varepsilon>0\)，选 \(m\) 使 \(n\ge m\) 时
\(|S_n-A|<\varepsilon\)。则（15）右边的尾和绝对值至多为
\(\varepsilon V\)，与 \(h\) 无关。对 \(n<m\) 的有限项，\(h\to0\) 时
\(w_n,w_{n+1}\to1\)，故每一项以及首项 \(A(w_1-1)\)
均趋于零。取上极限再令
\(\varepsilon\downarrow0\)，即得（13）。这里允许原级数只条件收敛。

取 \(K(t)=(\sin t/t)^2\) 并规定 \(K(0)=1\)。它在零点附近的导数为
\(O(t)\)，在无穷远处有

\begin{equation}
K'(t)=\frac{\sin(2t)}{t^2}-\frac{2\sin^2 t}{t^3}=O(t^{-2}).
\tag{16}
\end{equation}

所以 \(\int_0^\infty|K'(t)|\,dt<\infty\)，且
\(K(t)\to0\)，这就还原了原引理。另取 \(K(t)=e^{-t}\) 得 Abel 权；取
\(K(t)=\max(1-t,0)\) 及 \(h=1/N\)，所得加权和恰为
\((S_0+\cdots+S_{N-1})/N\)，即部分和的算术平均。两例都由同一个尾部估计推出。

\par\medskip\noindent\phantomsection\addcontentsline{toc}{section}{方法总结}\textbf{方法总结}\quad B317---B320 的共同问题是怎样在较大的函数空间里定义 Fourier
变换。先对逼近列的差使用已有估计，以完备性定义极限，再用同一列的两种收敛方式识别两种定义。B321
则要分别检查乘积所属的空间和卷积逐点存在的条件。B322
将相同的思路用于数项级数，先分部求和，再以核的总变差控制所有小参数下的尾和。

\end{document}
